\section{Introduction}
\label{sec:intro}

A difference equation is \emph{globally periodic} if all of its orbits are
periodic with a common period. The celebrated example is the Lyness
equation
\begin{equation}
x_{n+2} = \frac{1+x_{n+1}}{x_n},
\end{equation}
every positive solution of which repeats with period $5$
\citep{Lyness1942}. Global periodicity is a fragile property, and
equations that have it are rare. In a survey of open problems in
low-dimensional dynamics, \citet{Gasull2021} records five globally periodic
rational difference equations with nonnegative coefficients,
\begin{equation}
x_{n+1}=x_n,\quad
x_{n+1}=\frac{1}{x_n},\quad
x_{n+2}=\frac{x_{n+1}}{x_n},\quad
x_{n+2}=\frac{1+x_{n+1}}{x_n},\quad
x_{n+3}=\frac{1+x_{n+1}+x_{n+2}}{x_n},
\label{eq:fivefamilies}
\end{equation}
of periods $1$, $2$, $6$, $5$ and $8$ respectively. The $8$-cycle of the
last was found by H.~Todd \citep{Lyness1945,KocicLadasRodrigues1993}. New examples can be
made from these in two trivial ways: by \emph{rescaling}, $y_n=\alpha x_n$
with $\alpha>0$, and by \emph{unfolding}: multiplying every index in an equation of
order $k$ by an integer $\ell\ge1$ gives an equation of order $\ell k$
whose solutions are the interleavings of $\ell$ solutions of the original
one, and which is therefore again globally periodic (Gasull uses the word
differently; see \autoref{sec:reduction}). Gasull's Problem~11
asks whether there are any others. More precisely, consider the equations of order
$k\ge1$,
\begin{equation}
x_{n+k}=\frac{A_0+A_1x_n+A_2x_{n+1}+\cdots+A_kx_{n+k-1}}
             {B_0+B_1x_n+B_2x_{n+1}+\cdots+B_kx_{n+k-1}},
\label{eq:family}
\end{equation}
with $A_i,B_i\ge 0$, $\sum_i A_i>0$ and $\sum_i B_i>0$, acting on positive
initial data. (Gasull also assumes $A_1^2+B_1^2\ne0$, so that the equation
genuinely involves $x_n$; I do not need this, because global periodicity
forces it, as the first step of the proof shows. He also calls an equation
$p$-periodic only when $p$ is the least common period, which makes no
difference to the question.) Is every globally periodic equation of the form
\eqref{eq:family} obtained from one of \eqref{eq:fivefamilies} by
rescaling and unfolding?

Here I show that the answer is yes at every order.

\begin{theorem}
\label{thm:main}
Let \eqref{eq:family} be globally periodic, that is, suppose there is a
$p\ge1$ such that every positive solution satisfies $x_{n+p}=x_n$. Then,
up to multiplying the numerator and denominator by a common positive
constant, \eqref{eq:family} is one of the following equations, written from
the second on with the denominator cleared:
\begin{enumerate}
\item the shift $x_{n+k}=x_n$;
\item $x_{n+k}\,x_n=a$, with $a>0$;
\item $k=2m$ and $x_{n+2m}\,x_n=b\,x_{n+m}$, with $b>0$;
\item $k=2m$ and $x_{n+2m}\,x_n=b^2+b\,x_{n+m}$, with $b>0$;
\item $k=3m$ and $x_{n+3m}\,x_n=b^2+b\,(x_{n+m}+x_{n+2m})$, with $b>0$.
\end{enumerate}
These are, respectively, the $k$-fold unfoldings of the first two
equations of \eqref{eq:fivefamilies} and the $m$-fold unfoldings of the
last three, after rescaling; so every globally periodic equation of the
form \eqref{eq:family} is equivalent to one of the five equations of
\eqref{eq:fivefamilies}.
\end{theorem}

Conversely, each equation in the list is globally periodic, with periods
$k$, $2k$, $6m$, $5m$ and $8m$ respectively, since it is an unfolding of a
rescaling of one of \eqref{eq:fivefamilies}; so \autoref{thm:main} is a
complete classification.

The requirement of a \emph{common} period loses nothing: if every orbit in
$(0,\infty)^k$ is periodic, the state map is a continuous bijection, hence a
homeomorphism by invariance of domain, and a pointwise periodic homeomorphism
of a connected manifold is periodic \citep{Montgomery1937}; often it can even
be linearised globally \citep{CimaGasullManosas2012}. (Discontinuous maps can
be pointwise periodic without being periodic
\citep{CimaGasullManosaManosas2022}.) This remark is not used in the proof.

The question has a long history. Essentially the same question was asked by
\citet[Research problem~1.24]{GrahamKnuthPatashnik1989}; the generalisation that
Graham had in mind for his Problem~1343 in \emph{Mathematics Magazine} asks for
equations of much the same form \citep[Remark~1]{KocicLadasRodrigues1993}. A version of the
question is posed by \citet[p.~161]{GroveLadas2005}; at order two it is Conjecture~6.5.6 of
\citet{KulenovicLadas2000}. After partial results of
\citet{KocicLadasRodrigues1993}, order two was settled by
\citet{BarbeauGelbordTanny1995}, and \citet{CsornyeiLaczkovich2001} settled,
even for complex coefficients, orders two and three and the reciprocal
equations $x_{n+k}x_n=c_0+\sum_{j=1}^{k-1}c_jx_{n+j}$ in which every $c_j$
is nonzero (the full support). \citet{CimaGasullManosas2004} reduced the
problem to a reciprocal family with palindromic coefficients, and showed,
using a continuum of $2$-periodic points at odd order (\autoref{sec:oddk}),
that the answer is yes for $k\in\{1,2,3,4,5,7,9,11\}$. The general case has
remained open. Different approaches to
the global periodicity problem are surveyed by
\citet{CimaGasullManosaManosas2016}.
Over the complex numbers and with arbitrary coefficients, periodic linear
fractional recurrences have been studied from the point of view of
birational geometry by \citet{BedfordKim2006,BedfordKim2014} at orders $2$
and $3$, and \citet{BedfordKim2011} show that for every $k$ there are
such recurrences of order $k$ and period $4k$;
positivity of the coefficients, which is essential here, plays no role
there.

The proof has two quite different halves. The first is analytic and
short. A surjectivity argument, followed by a Lyapunov function and the
derivative at a fixed point, shows that a globally periodic equation is
either the shift or of the \emph{reciprocal} form
\begin{equation}
x_{n+k}\,x_n = c_0 + \sum_{j\in S} c_j\,x_{n+j},
\qquad c_0\ge0,\quad c_j > 0 \text{ for } j \in S \subseteq \{1,\dots,k-1\},
\end{equation}
and linearising at the interior equilibrium shows that the coefficients
are palindromic, $c_j=c_{k-j}$. Throughout, $\theta=1$ if $c_0>0$ and
$\theta=0$ if $c_0=0$; most of the arguments treat both values at once. Unfolding then lets me assume that the
\emph{support} $S$ is primitive, $\gcd(S\cup\{k\})=1$. When $S$ has at
most two elements, the problem reduces to equations of order at most
three, which are settled with the min-plus recursion described below
(and, at order three with $c_0=0$, with the linearisation), and to a spectral statement about the
polynomial $\lambda^k+1-\eps(\lambda^j+\lambda^{k-j})$, which is settled
with Kronecker's theorem on algebraic integers whose conjugates all lie
on the unit circle.

The second half, which is most of the paper, deals with supports of three
or more elements, and it is combinatorial. The key observation is that
global periodicity survives degeneration. If the initial data are taken
to be powers $\delta^{e_i}$ of a small parameter $\delta$, the exponents of
the solution follow a piecewise-linear (min-plus) recursion that depends
only on the support and on $\theta$. For the Lyness equation, for example,
$x_n\approx\delta^{e_n}$ turns $x_{n+2}x_n=1+x_{n+1}$ into
$e_{n+2}+e_n=\min(0,e_{n+1})$, and every integer solution of this recursion
has period $5$, just as the solutions of the equation do. In general a
globally periodic equation forces its min-plus recursion to be periodic.
Because all the coefficients are positive, there is no cancellation: each
term of the solution behaves like a positive multiple of a power of
$\delta$, and these leading coefficients follow a recursion of their own,
which must also be periodic; this is what settles orders two and three when
$c_0>0$. The passage to a piecewise-linear recursion is the limiting
procedure of \citet{TokihiroTTMS1996}, now called ultradiscretisation and
familiar from integrable systems
\citep{Nobe2008,InoueTakenawa2008,InoueKunibaTakagi2012}; for the full
support it was already used by \citet{CsornyeiLaczkovich2001}. Tropical
versions also appear in the theory of cluster algebras, where
\citet{FominZelevinsky2003} use them to prove the periodicity of
Zamolodchikov's $Y$-systems, of which the Lyness equation is the case of type
$A_2$; here the information flows the other way, from the equation to its
tropical version.

For larger supports, global periodicity of an equation is refuted, for
\emph{every} choice of positive coefficients at once, by exhibiting an
unbounded solution of the \emph{tangent recursion}, the linearisation of the
min-plus recursion along a periodic solution. The idea is to perturb a
periodic orbit of the min-plus recursion slightly. Only the terms that attain
the minimum see the perturbation, so it follows a recursion of the same kind,
and if the perturbation can grow without bound, the min-plus recursion, and
with it the equation, cannot be periodic. For the supports with at least three
elements, the periodic orbits used have, except at order $6$, a minimum that is
always negative, so that the constant term never attains it, and both values of
$\theta$ are refuted at once.

I construct the unbounded tangent solutions explicitly, case by case
(\autoref{tab:intro-cases}). Apart from the full support and most supports
at odd order, which have tangent solutions of their own, the workhorse is the
\emph{$\sigma$-bump}: a linearly growing tangent solution supported on two
positions in every block of $k$, whose existence reduces to a purely
combinatorial condition on a set of residues modulo $2k$
(\autoref{def:bump}). The hardest case is even order with every element of
$S$ odd, where the residues are chosen as a half-arc
$\{y: \nu y \bmod 2k \in [0,k)\}$ for a suitable odd multiplier $\nu$, and
what is needed is a lemma about three-term arithmetic progressions under
unit multipliers modulo $2k$, proved with the geometry of numbers in a plane
lattice of determinant $2k$ and Jacobsthal's function.

Several steps finish with a finite computation, for instance at small
orders where the general arguments do not yet apply; each is stated
precisely in the text and listed in \autoref{tab:intro-cases}. Every finite
computation is carried out in exact integer arithmetic, and the whole proof,
including these computations, has been formalised and kernel-checked in the
\textsc{Lean} proof assistant (\autoref{app:lean}).

As far as I know, the following ingredients are new: the tangent recursion
and the two-scale pinning lemma (\autoref{lem:trop-pinning}), which turn an
unbounded solution of the linearised min-plus recursion into a proof that the
equation is not periodic; the use of such tangent solutions to refute global
periodicity for every choice of coefficients at once; the $\sigma$-bump and
its carriers (the parity, band and half-arc carriers); and the theorem on
three-term progressions under unit multipliers (\autoref{thm:lemmaA}).
Ultradiscretisation itself is not new, nor is the use of the leading
coefficients at orders two and three, which goes back to
\citet{CsornyeiLaczkovich2001}. I am not aware of an earlier use of the derivative of a min-plus map along a
periodic orbit to show that the map is not periodic (``linearised'' means
something different in the ultradiscrete literature, namely straight-line
motion on a tropical Jacobian). A by-product is a purely
tropical statement: for every palindromic support with at least three
elements, and either $\theta$, no power of the min-plus map is the identity
(\autoref{cor:trop-nonperiodic}).

\autoref{sec:reduction} reduces the problem to primitive palindromic
supports of the reciprocal form and \autoref{sec:tropical} develops the
min-plus machinery and the $\sigma$-bump; Sections~\ref{sec:small}--\ref{sec:lemmaA}
treat the cases of \autoref{tab:intro-cases}, \autoref{sec:assembly} puts
them together, and the appendices hold the finite range of
\autoref{thm:lemmaA} and the bounds on Jacobsthal's function, the most
mechanical case checks, and a description of the formalisation. A reader who wants only the shape of the argument can
read Sections~\ref{sec:reduction} and~\ref{sec:tropical}, the statements of
the main theorems of the later sections, and \autoref{sec:assembly}.

\begin{table}
\centering
\footnotesize
\renewcommand{\arraystretch}{1.15}
\begin{tabular}{@{}>{\raggedright\arraybackslash}p{0.15\textwidth}%
>{\raggedright\arraybackslash}p{0.19\textwidth}%
>{\raggedright\arraybackslash}p{0.27\textwidth}%
>{\raggedright\arraybackslash}p{0.23\textwidth}%
>{\raggedright\arraybackslash}p{0.05\textwidth}@{}}
case & mechanism & results used & finite computations & \S\\\hline
every support & reduction to \eqref{eq:familya}, palindromy, folding &
Theorems~\ref{thm:red-dichotomy}, \ref{thm:red-shift},
\ref{thm:reduction}, \ref{thm:spectral}, \ref{thm:palindromy};
\autoref{lem:unfold} & none & \ref{sec:reduction}\\
$k=2$, $S=\{1\}$ & leading coefficients & \autoref{thm:order2}, from
Theorems~\ref{thm:leading} and \ref{thm:P} & none & \ref{sec:small}\\
$k=3$, $S=\{1,2\}$ & leading coefficients; for $c_0=0$, the linearisation &
\autoref{thm:order3}, from Theorems~\ref{thm:leading}, \ref{thm:P} and (for
$c_0=0$) \ref{thm:spectral} & none & \ref{sec:small}\\
two-term supports $\{j,k-j\}$ & the linearisation and Kronecker's theorem &
\autoref{thm:twoterm}, from \autoref{thm:spectral},
Lemmas~\ref{lem:small-newton} and \ref{lem:small-bound} & none &
\ref{sec:small}\\
full support, $k\ge4$ & a closed-form tangent solution on the carrier
$Q=\{1\}$ & \autoref{thm:fullsupport}, from Lemmas~\ref{lem:fs-carrier}
and \ref{lem:fs-ray} and \autoref{thm:tangent} & none &
\ref{sec:fullsupport}\\
odd $k$, $X$ not a dilated interval & rays along the alternating orbit &
\autoref{thm:odd-nondil}, from \autoref{prop:odd-second},
\autoref{lem:odd-ray}, the lemmas of \autoref{sec:odd-rigid} and
\autoref{thm:tangent} & none & \ref{sec:odd-nondil}\\
odd $k$, $X$ a dilated interval & a $\sigma$-bump &
Theorems~\ref{thm:odd-rowcarrier} and \ref{thm:sigmabump},
\autoref{lem:odd-fold} & none & \ref{sec:odd-dil}\\
even $k$, some element of $S$ even & $\sigma$-bumps on parity carriers, or
on two explicit carriers & \autoref{thm:notallodd}, from
Theorems~\ref{thm:nao-evenmin} and \ref{thm:nao-struct} and
\autoref{cor:nao-bothfailing}; \autoref{thm:sigmabump} & none &
\ref{sec:notallodd}\\
even $k$, all odd, $1\in S$ & $\sigma$-bumps on small explicit carriers &
\autoref{prop:allodd-one}; \autoref{thm:sigmabump} & none &
\ref{sec:allodd}\\
even $k$, all odd, $1\notin S$, $|L|=1$ & $\sigma$-bumps on half-arc
carriers & \autoref{thm:allodd-threeall} (all but five supports);
\autoref{thm:sigmabump} & $k\le22$: five carriers of
\autoref{tab:allodd-base} & \ref{sec:allodd}\\
even $k$, all odd, $1\notin S$, $|L|\ge2$ & $\sigma$-bumps on half-arc
carriers, with \autoref{thm:lemmaA} & \autoref{prop:allodd-e3bump},
\autoref{thm:allodd-e3all}, \autoref{thm:lemmaA}; \autoref{thm:sigmabump} &
\autoref{lem:allodd-pf}; $k\le90$ in \autoref{thm:allodd-e3all}; twelve carriers of
\autoref{tab:allodd-base}; for \autoref{thm:lemmaA}, the chord table (C1),
the chains of \autoref{tab:lemmaA-chains}, the certificate (C2)--(C3) at the
$127$ orders $92\le k\le2730$ with three or more odd prime factors, and the
checks in
Lemmas~\ref{lem:lemmaA-pattern} and \ref{lem:lemmaA-crude} &
\ref{sec:allodd}, \ref{sec:lemmaA}\\
$k=6$, $S=\{1,3,5\}$ & a tangent certificate &
Propositions~\ref{prop:order6cert} and \ref{prop:trop-cert};
\autoref{thm:tangent} & the certificate of \autoref{prop:order6cert} &
\ref{sec:tropical}
\end{tabular}
\caption{The cases of the proof, the mechanism and results each uses, and
the steps that rest on a finite computation (``none'' means that the case is
proved by hand for all parameters). Every case uses the reduction of the
first row. Except in the rows for $|S|\le2$, $|S|\ge3$, and $S$ is not the
full support unless stated; $\theta$ is arbitrary throughout. At odd order
$k$, $X$ is the set of halved even elements of $S$ (\autoref{lem:odd-coords});
$L=\{s\in S:2s<k\}$.}
\label{tab:intro-cases}
\end{table}
